\documentclass{tufte-handout}

%\geometry{showframe}% for debugging purposes -- displays the margins

\usepackage[spanish, es-tabla]{babel}
\usepackage{amsmath}
\usepackage{chemfig}


% Set up the images/graphics package
\usepackage{gensymb}
\usepackage{physics}
\usepackage{multicol}

\usepackage{graphicx}
\setkeys{Gin}{width=\linewidth,totalheight=\textheight,keepaspectratio}
\graphicspath{{graphics/}}

\title[Prácticas Química Física II: UV-Vis]{
Prácticas Laboratorio 2: Espectroscopía UV-Visible}
%\author{David De Sancho}
\date{}  % if the \date{} command is left out, the current date will be used

% The following package makes prettier tables.  We're all about the bling!
\usepackage{booktabs}

% The units package provides nice, non-stacked fractions and better spacing
% for units.
\usepackage{units}

% The fancyvrb package lets us customize the formatting of verbatim
% environments.  We use a slightly smaller font.
\usepackage{fancyvrb}
\fvset{fontsize=\normalsize}

% Small sections of multiple columns
\usepackage{multicol}

% Provides paragraphs of dummy text
\usepackage{lipsum}

% These commands are used to pretty-print LaTeX commands
\newcommand{\doccmd}[1]{\texttt{\textbackslash#1}}% command name -- adds backslash automatically
\newcommand{\docopt}[1]{\ensuremath{\langle}\textrm{\textit{#1}}\ensuremath{\rangle}}% optional command argument
\newcommand{\docarg}[1]{\textrm{\textit{#1}}}% (required) command argument
\newenvironment{docspec}{\begin{quote}\noindent}{\end{quote}}% command specification environment
\newcommand{\docenv}[1]{\textsf{#1}}% environment name
\newcommand{\docpkg}[1]{\texttt{#1}}% package name
\newcommand{\doccls}[1]{\texttt{#1}}% document class name
\newcommand{\docclsopt}[1]{\texttt{#1}}% document class option name

\begin{document}

\maketitle% this prints the handout title, author, and date

\begin{abstract}
En la segunda práctica utilizaremos otro tipo de espectroscopía,
la que usa radiación ultravioleta-visible (correspondiente a 
longitudes de onda entre 170 y 780 nm). Este tipo de radiación,
más energética que la infrarroja, afecta a los niveles electrónicos 
de las moléculas. Aunque en este caso no podemos obtener expresiones
analíticas compactas para las energías correspondientes a las transiciones, sí que podemos aprender cómo el entorno de un enlace
químico afecta la energía de los orbitales. 
\end{abstract}

\section{Introducción}
La energía de un cuanto de radiación con energía $h\nu$ puede
resultar en la excitación de niveles energéticos electrónicos.
La intensidad de absorción se expresa mediante los 
\textbf{coeficientes de absorción molar}, $\varepsilon$. 
Para una solución, este coeficiente está relacionado con la 
\textbf{absorbancia} a través de la relación
\begin{equation}
    \varepsilon = \frac{A}{cl},
\end{equation}
donde $c$ es la concentración en moles por litro y $l$ es la
longitud de la celda en cm. 
Un parámetro característico de una 
molécula es la longitud de onda en la que encontramos
el máximo de absorción, $\lambda_{max}$.

\subsection{Interpretación de los espectros}
En el caso de moléculas poliatómicas, podemos relacionar 
las transiciones electrónicas observadas a excitaciones de
electrones entre distintos tipos de orbitales moleculares.
Las partes de las moléculas responsables de transiciones 
que podemos observar en esta región del espectro 
electromagnético se llaman \textbf{cromóforos}. En la Tabla 
\ref{tab:chrom} mostramos algunos cromóforos de moléculas
orgánicas a las que a menudo se aplica este tipo de
espectroscopía, los tipos de orbitales implicados y 
las longitudes de onda aproximadas a las que se observan. 
Típicamente, observamos transiciones entre el orbital molecular
ocupado de más alta energía (HOMO) y el orbital desocupado de más
baja energía (LUMO). En función de los diferentes tipos de orbitales
moleculares implicados en las transiciones, las transiciones
serán más o menos energéticas. 

\begin{table}
\centering
    \begin{tabular}{|c|c|c|c|}
    \hline
     Cromóforo & Transición &  $\lambda_{max}$ (nm) & Intensidad \\
     \hline
     \chemfig{C=C} & $\pi\rightarrow\pi^\star$ & 190 & media\\
    \hline  
     \chemfig{C=O} & $n\rightarrow\pi^\star$& 300 & débil\\
      & $n\rightarrow\sigma^\star$& 190 & intensa\\
    \hline  
     \chemfig{C=C-C=C} & $\pi\rightarrow\pi^\star$ & 210 & intensa \\
    \hline  
     \chemfig{C=C-C=O} & $n\rightarrow\pi^\star$ & 330 & débil \\
     \chemfig{C=C-C=O} & $\pi\rightarrow\pi^\star$ & 210 & intensa \\
      & $\pi\rightarrow\pi^\star$ & 210 & intensa \\
    \hline  
     & $\pi\rightarrow\pi^\star$ & 180 & muy intensa \\
& $\pi\rightarrow\pi^\star$ & 200 & intensa \\
      \chemfig{*6(=-=-=-)} & $\pi\rightarrow\pi^\star$ & 255 & débil\\
\hline
\end{tabular}
\caption{Algunos cromóforos de moléculas orgánicas, tipos de
electrones implicados, las longitudes de onda aproximadas e intensidad de las bandas correspondientes.}\label{tab:chrom}
\end{table}

En el caso de los enlaces \chemfig{C=C} observamos la excitación
de electrones $\pi$ que llegan a orbitales $\pi^\star$ 
antienlazantes. A estas transiciones se las llama
$\pi\rightarrow\pi^\star$ y su energía corresponde al espectro UV. 
Cuando el doble enlace se encuentra en una cadena conjugada, 
las energías de los orbitales moleculares se acercan entre sí
y pueden llegar al rango del visible. En el caso del enlace
\chemfig{C=O} nos encontramos un electrón no enlazante $n$
que puede excitarse para llegar a un orbital vacío $\pi^\star$
antienlazante (transición $n\rightarrow\pi^\star$). Estas
transiciones en realidad están prohibidas por simetría y por
tanto suelen ser bastante débiles.

\section{Objetivos}
En esta práctica mediremos la longitud de onda de absorción
máxima para la acetona en diferentes disolventes y de una
serie de derivados de la acetona con diferentes sustituyentes 
en ciclohexano. Los disolventes formarán \textbf{enlaces de hidrógeno}
con la molécula en distinto grado.
Los sustituyentes modificarán por \textbf{efecto 
inductivo} o por \textbf{efecto de resonancia} la longitud de onda
de absorción del enlace.  

\section{Protocolo}
\subsection{Material}
    \begin{multicols}{2}
\begin{itemize}
\item Celda de cuarzo para UV
\item Peras aspirapipetas
\item Matraces aforados de 50 ml
\item Pipetas de 1 ml
\end{itemize}
    \end{multicols}

\subsection{Reactivos}
    \begin{multicols}{2}
\begin{itemize}
\item Ciclohexano
\item Cloroformo
\item Acetonitrilo
\item Etanol
\item Agua
\item Acetona
\item Cloruro de acetilo
\item Ácido acético
\item Metiletilcetona
\item Acetofenona
\end{itemize}
    \end{multicols}

\subsection{Procedimiento experimental}
\begin{enumerate}
    \item Se preparan disoluciones (50 ml) 0,04 M de acetona en cada uno de los siguientes disolventes: ciclohexano, cloroformo, acetonitrilo, etanol y agua.
    \item Se registran los espectros UV de las disoluciones de la acetona, con los disolventes como referencia, y se mide el valor de $\lambda_{max}$ para la banda correspondiente a la transición $n\rightarrow\pi^\star$ en cada disolución, en la región comprendida entre 200 y 350 nm.
    \item Se preparan disoluciones aproximadamente $0.04$ M en ciclohexano de cada uno de los siguientes compuestos carbonílicos: metiletilcetona, acetona, ácido acético, cloruro de acetilo y acetofenona.
    \item Se registran los espectros UV de las disoluciones de los compuestos citados en el apartado anterior, con el disolvente (ciclohexano) como referencia, y se mide el valor de $\lambda_{max}$ para la banda correspondiente a la transición $n\rightarrow\pi^\star$ en cada disolución, en la región comprendida entre 200 y 350 nm.
\end{enumerate}

\subsection{Cálculos teóricos}
Los métodos mecanocuánticos nos permiten identificar los
valores de $\lambda_{max}$ de moléculas sencillas en distintos
disolventes. Para ello, utilizaremos los programas Gaussian y
Gaussview. Además, visualizaremos los orbitales implicados en
las transiciones electrónicas y mediremos las energías 
asociadas a las excitaciones electrónicas.

\section{Resultados y discusión}
\begin{itemize}
    \item Expresar en forma de tabla los valores de $\lambda_{max}$ obtenidos en los diferentes disolventes.
\item Interpretar las variaciones observadas en función de la naturaleza de los disolventes.
\item Expresar en forma de tabla los valores de $\lambda_{max}$ obtenidos para los diferentes compuestos
carbonílicos.
\item Correlacionar los valores de las longitudes de onda con el carácter dador o aceptor de los
sustituyentes unidos al grupo acetilo, y compararlo con el dato bibliográfico del
acetaldehído.
\end{itemize}
\end{document}